% Review-only v10 prefix closure construction, 9 October 2026.
% Candidate TODO at the passage asserting Kpt is a finitely
% branching tree, immediately after the definition of Kpt.
% Consolidate with preceding first-full-prefix/RawTypeMeet notes.
% No mathematical manuscript prose is automatically modified.
%
\todo[inline]{\textbf{v10 validation (prefix closure of partial types).}
To justify that every initial-socle restriction of a node
$\operatorname{type}_{A^+}(v)$ is itself an admissible
partial type, the type vertex cannot simply be placed at
the new socle end: the last required $E$-pair would be
consecutive, contradicting Definition~6.28(1).
Instead, retain the first $d$ ordinary vertices, add
one neutral filler vertex at position $d$, and place
the distinguished original vertex at position $d+1$.
Give it precisely the $E$-socle $\{0,\ldots,d-1\}$
and copy the original $L$-type over that socle.
The $E$-part of this construction has a unique free
cut at $d$, satisfies spacing and downward closure,
and preserves every old and crossing $E$-bit,
as checked in \texttt{PrefixReplicaE.lean}.
Still verify the induced $L$-relations, the third
partial-structure axiom, the filler singleton's
allowed type and preservation of forbidden-freeness
(the latter uses irreducibility). Only then is the
actual $K^\mathcal F_{pt}$ family proved prefix-closed.}
%
% Proof endpoints to check at exact-head green CI:
%  prefixReplicaE_old
%  prefixReplicaE_last_iff
%  prefixReplicaE_spaced
%  prefixReplicaE_downward
%  prefixReplicaE_last_freeCut
%  prefixReplicaE_last_freeLevel
%  prefixReplicaE_preserves_type_socle
%  prefixReplicaE_preserves_old_socle
